\section{Introduction}

Rank polymorphism is a language feature that permits the definition of
functions operating on arrays of arbitrary rank, where the rank of an array is
the number of its dimensions (axes). As an illustration, consider a
function that adds two numeric arrays element-wise; the same definition applies
to 0-dimensional, 1-dimensional, $\ldots$ arrays. A common justification for
supporting rank polymorphism is its facilitation of generic programming: rather
than providing a separate definition for each array rank, a single definition
suffices for all ranks. This reduces the likelihood of errors, curtails
verbosity, and---because compositions of rank-polymorphic operations yield more
general specifications---encourages programmers to contemplate corner cases that
might otherwise be overlooked~\citep{10.1145/1283920.1283935}. Arguably, this
is a principal reason why specifications in the APL~\citep{Iverson:1962:APL}
family are exceptionally terse.

Although this interpretation of rank polymorphism dates back to the 1960s, the
present pearl offers a fresh perspective on the core idea, showing that,
when rank polymorphism is made to work harder, it yields considerably richer
benefits.

First, rank polymorphism gives rise to a notion of algorithmic
equivalence. For a given rank-polymorphic function one can often define an
equivalence relation on array shapes such that reshaping an array within the
same equivalence class does not affect the computed values. A simple example is
the sum of all array elements---the resulting value remains unchanged regardless
of the shape assigned to the input array (provided that reshapes neither add
nor remove elements).

Second, if array shapes are endowed with a richer structure than a mere
sequence of dimensions, rank polymorphism becomes a tool for performance
engineering. As \citet{10.5555/280586} demonstrated, the choice of data structures determines
algorithmic complexity in functional languages. Analogously, we show that the
choice of array shape influences the complexity of array-based algorithms.
Consider a rank-polymorphic function together with a shape equivalence relation
such that all inputs in the same class produce identical results. Assuming a
fixed evaluation strategy and runtime costs, reshapes within the equivalence
class dictate the actual execution time while preserving the computed values.
In other words, a rank-polymorphic specification defines \emph{what} is
computed, whereas the choice of shape within an equivalence class determines
\emph{how} the computation proceeds.  Consequently, finding an efficient
algorithm amounts to searching for a shape that minimises the runtime cost; in
this paper we compare candidate shapes by evaluating a cost model.  In
principle, such cost models can encode assumptions about the underlying
hardware, such as parallelism, memory hierarchy or cache sizes; here we only
model the number of operations and parallelism.  Other aspects, such as the
layout of arrays in memory or the sparsity of the data, are determined by the
choice of array representation rather than by shapes, and our cost model does
not account for them.

The principal technical contribution of this work is a precise formulation of
the above ideas via a framework in Agda, which enables us to define:
%
(i) rank-polymorphic arrays (Section~\ref{sec:arrays});
(ii) equivalence classes of shapes (Section~\ref{sec:reshaping}); and
(iii) runtime-cost models (Section~\ref{sec:runtime}).

Using the scan algorithm as a running example, we (i) encode a rank-polymorphic
version of scan; (ii) prove that all scans belonging to a chosen
shape-equivalence class compute identical results; (iii) devise a cost model
for an idealised multicore processor; and (iv) demonstrate the runtime impact
of reshaping.

When we introduce rank-polymorphic arrays, we deliberately avoid fixing the
representation of arrays to any particular data structure. We achieve this by
specifying an array theory that stipulates assumptions about shapes and
indices, whilst requiring the concrete array datatype to be a representable
functor. Our running example and accompanying equivalence proof are presented
without committing to a specific model, meaning that they can be instantiated
for any admissible representation. Thus, we may (a) define array
representations and basic combinators in Agda; (b) perform performance
engineering to select optimal array shapes for a given representation; and (c)
subsequently instantiate our algorithm for that representation, preserving the
established correctness guarantees.  In the remainder of the paper we guide
the reader through the rationale behind our design choices and the technical
details of each component of the framework.  This paper is an Agda
script, available online~\citep{paper-src}; it has been checked with
Agda 2.8.0~\citep{agda-2-8-0} and the Agda standard library
version 2.3~\citep{agda-stdlib-2-3}.
